\chapter{Hamiltonian symmetry and continuous cochains}
\label{chap:hamiltonian-symmetry-cochains}

An infinitesimal change of two coordinates can preserve their area
form while depending on arbitrarily many Taylor coefficients.
Translations are among these changes. Their action lowers Taylor
order, so truncating the coordinates can destroy a relation between
symmetries. To retain the full action, the relations must be formed
before truncation. This leads to a complex of continuous cochains
whose coefficients have different, individually finite sets of
symmetry parameters.

\section{Area-preserving changes of formal coordinates}

Fix a field $k$ of characteristic zero, with the discrete topology.
Write
\[
 R=k[[x,y]],\qquad \mathfrak m=(x,y).
\]
An element of $R$ is a family of coefficients $f_{ab}\in k$, written
$\sum_{a,b\geq0}f_{ab}x^ay^b$. Products are defined coefficientwise.
Each coefficient of a product is a finite sum. The neighborhoods of
zero are the subspaces $\mathfrak m^N$ of series whose terms have
total degree at least $N$. Thus convergence means eventual agreement
of every finite set of coefficients. This is the algebraic formal
topology; in particular the coefficient field remains discrete when
$k=\mathbb C$.

Formal partial differentiation is defined term by term. It is
continuous because the coefficients through degree $N$ of a derivative
depend on coefficients through degree $N+1$ of the original series.
A continuous derivation of $R$ is a continuous $k$-linear map $V$
satisfying $V(fg)=V(f)g+fV(g)$. It has the unique form
\[
 V=a\partial_x+b\partial_y,\qquad a=V(x),\quad b=V(y).
\]
The formula follows first on monomials by the product rule and then on
series by continuity.

Introduce a variable $\varepsilon$ with $\varepsilon^2=0$. The change
$(x,y)\mapsto(x+\varepsilon a,y+\varepsilon b)$ transforms the
alternating area form by
\[
 d(x+\varepsilon a)\wedge d(y+\varepsilon b)
 =\bigl(1+\varepsilon(\partial_xa+\partial_yb)\bigr)
       dx\wedge dy.
\]
Here $df=(\partial_xf)dx+(\partial_yf)dy$, and the wedge product
satisfies $dx\wedge dx=dy\wedge dy=0$ and
$dy\wedge dx=-dx\wedge dy$. Area preservation is therefore the equation
$\partial_xa+\partial_yb=0$.

Every solution is generated by a formal function. Indeed, define
\[
 f(x,y)=\int_0^x b(s,y)\,ds-\int_0^y a(0,t)\,dt.
\]
These are formal antiderivatives, obtained by dividing a monomial by
its new positive exponent. Characteristic zero permits these divisions.
The area equation gives $\partial_xf=b$ and $\partial_yf=-a$.
Thus $V=H_f$, where
\[
 H_f=-(\partial_yf)\partial_x+(\partial_xf)\partial_y,
 \qquad
 H_f(g)=\{f,g\}:=\partial_xf\,\partial_yg
                        -\partial_yf\,\partial_xg.
\]
Such an $f$ is called a Hamiltonian for $V$. Its ambiguity is a constant:
both partial derivatives of a series vanish precisely when all its
nonconstant coefficients vanish.

The bracket of two derivations is their commutator,
$[V,W]=V\circ W-W\circ V$. Expanding the two compositions gives
$[H_f,H_g]=H_{\{f,g\}}$. Commutators obey the Jacobi identity because
each ordered composition of three operators cancels with its opposite
occurrence. It follows that the space
\[
 \mathfrak h=R/k
\]
is a Lie algebra, with bracket induced by $\{f,g\}$. A Lie algebra here
means a vector space with a bilinear antisymmetric bracket satisfying
$[[f,g],h]+[[g,h],f]+[[h,f],g]=0$. The map from $\mathfrak h$ to the
area-preserving continuous derivations is an isomorphism of Lie algebras.

Represent a class in $R/k$ by its series with zero constant coefficient.
Constants must still be removed \emph{after} forming its bracket.
For example,
\[
 \{x,y\}=1\quad\text{in }R,
 \qquad [x,y]=0\quad\text{in }\mathfrak h,
 \qquad H_x=\partial_y,\quad H_y=-\partial_x.
\]
The translations commute as vector fields. The functions $1,x,y$
retain a central bracket before the quotient. This difference between
Hamiltonians and their vector fields is part of the construction.

Let $H_d$ be the finite-dimensional space of homogeneous polynomials
of degree $d$. Then
\[
 \mathfrak h=\prod_{d\geq1}H_d.
\]
For positive $a,b$, taking brackets of terms of orders $a,b$ gives terms
of order at least $a+b-2$, modulo constants. More precisely, a
nonconstant output of degree at most $N$ uses inputs only through degree
$N+1$. This proves joint continuity of the bracket in the stated topology.

\section{Why a Taylor window is not a Lie quotient}

The first few brackets already distinguish the full symmetry from
symmetry fixing the origin:
\[
 [x,y^2]=2y,\qquad [x,y^3]=3y^2,\qquad
 [x^2,y^2]=4xy,\qquad [x^3,y^3]=9x^2y^2.
\]
The coefficient projection through degree $N$ kills $y^{N+1}$ but
does not kill its bracket with $x$:
\begin{equation}\label{eq:ham-core-translation-window}
 [x,y^{N+1}]=(N+1)y^N.
\end{equation}
Its kernel is therefore not a Lie ideal. A Lie ideal $J$ is a vector
subspace such that $[\mathfrak h,J]\subset J$; only such a subspace
permits a well-defined bracket on the quotient.

Truncating the bracket itself is worse in degree at least three.
Let $[f,g]_N$ mean the bracket followed by projection through degree $N$.
For $N\geq3$, the triple $(y,y^N,x^2y)$ has Jacobiator
\[
 [[y,y^N]_N,x^2y]_N+[[y^N,x^2y]_N,y]_N
       +[[x^2y,y]_N,y^N]_N=2Ny^N.
\]
The first inner bracket vanishes. The second equals $-2Nxy^N$
before projection and is discarded. The third inner bracket is $2xy$,
which contributes $2Ny^N$. Without the projection, the discarded term
would contribute $-2Ny^N$ and complete the cancellation.

\begin{proposition}\label{prop:ham-core-no-finite-quotient}
The Lie algebra $\mathfrak h$ has no proper open Lie ideal. Every
continuous Lie homomorphism from $\mathfrak h$ to a discrete Lie algebra
is zero. In particular, no inverse system of discrete finite-dimensional
Lie quotients represents $\mathfrak h$ with its given topology.
Moreover $[\mathfrak h,\mathfrak h]=\mathfrak h$.
\end{proposition}

\begin{proof}
If $J$ is open, it contains $\mathfrak m^N$ for some $N$.
For $f=\sum_{a+b\geq1}f_{ab}x^ay^b$, set
\[
 I_y^Nf=\sum_{a+b\geq1}\frac{b!}{(b+N)!}f_{ab}x^ay^{b+N}.
\]
It belongs to $\mathfrak m^{N+1}\subset J$ and satisfies
$(\operatorname{ad}_x)^NI_y^Nf=f$. If $J$ is a Lie ideal, repeated
brackets with $x$ keep this expression in $J$. Hence $J=\mathfrak h$.
The kernel of a continuous map to a discrete space is open, which
proves the homomorphism assertion. Taking one antiderivative gives
$f=[x,I_yf]$ and proves the last assertion.
\end{proof}

By contrast, $\mathfrak h_+=\mathfrak m^2$ is a Lie subalgebra.
Its elements generate vector fields vanishing at the origin. Each
$\mathfrak m^{N+1}$ is an ideal of $\mathfrak h_+$, since a bracket
with an element of order two cannot lower order. Thus
\[
 \mathfrak h_+\cong\varprojlim_N
       \mathfrak m^2/\mathfrak m^{N+1}.
\]
The cubic subalgebra has nilpotent finite quotients
$\mathfrak m^3/\mathfrak m^{N+1}$ for $N\geq2$: every bracket with a cubic
Hamiltonian raises the order by at least one. Neither limit includes
the translations. Even the quadratic part is already nonabelian:
\[
 e=x^2/2,\qquad f=-y^2/2,\qquad h=-xy,
 \qquad [h,e]=2e,\quad[h,f]=-2f,\quad[e,f]=h.
\]
These formulas define the usual three-dimensional Lie algebra
$\mathfrak{sl}_2$. They exhibit a genuine finite-dimensional symmetry
inside the origin-preserving theory.

\section{Coefficients and continuous symmetry parameters}

The failure of Lie quotients leaves the full bracket available.
Relations between its actions can be encoded by alternating functions
of symmetry parameters. Their continuity determines which infinite
expressions are permitted.

Index nonconstant monomials by
\[
 \mathcal I=\{(a,b)\in\mathbb N^2:a+b\geq1\},
 \qquad e_I=x^ay^b\quad(I=(a,b)).
\]
The continuous dual of $\mathfrak h$ is
\[
 D=\operatorname{Hom}_{\mathrm{cont}}(\mathfrak h,k)
    =\bigoplus_{I\in\mathcal I} k\epsilon^I,
 \qquad \epsilon^I(e_J)=\delta^I_J.
\]
Indeed, continuity to the discrete field makes the kernel of a
functional contain some $\mathfrak m^N$. Conversely, every functional
using finitely many coefficients is continuous. Give $D$ the discrete
topology.

Introduce commuting symbols $u_I$, one for each $I\in\mathcal I$.
A multi-index $\alpha=(\alpha_I)$ has nonnegative integer entries
and finite support. Write $u^\alpha=\prod_Iu_I^{\alpha_I}$ and
$w(\alpha)=\sum_I\alpha_I$. Define
\[
 A=\prod_\alpha k u^\alpha,
 \qquad O_f=\sum_I f_Iu_I
        \quad\left(f=\sum_I f_Ie_I\right).
\]
Multiplication sums over the finitely many decompositions
$\alpha=\beta+\gamma$. The topology is the product of the discrete
coefficient spaces. Finite coordinate polynomials are dense.
For $\lambda\in D$, the linear expression $O_f$ evaluates as
$\lambda(f)$. The general series in $A$ is formal; no convergence
at arbitrary numerical values of its variables is imposed.

Equivalently, $A$ is the inverse limit of the polynomial algebras on
the variables with $|I|\leq N$, truncated also at polynomial degree $N$.
Here $|I|=a+b$ is the degree of the Hamiltonian label. It is distinct
from the polynomial degree $w(\alpha)$ of a coefficient monomial.
The projections in this description are algebra maps.

The adjoint action defines an action on coefficients by
\[
 f\cdot O_g=O_{[f,g]},\qquad
 f\cdot(FG)=(f\cdot F)G+F(f\cdot G).
\]
It extends jointly continuously to $\mathfrak h\times A\to A$.
To check this, fix an output monomial. Each contributing term replaces
one input variable by an output variable $u_K$. Its structure coefficient
uses $[e_I,e_J]$ with output $e_K$, and therefore
$|I|+|J|=|K|+2$. Both input degrees are positive. Only finitely many
labels and coefficient monomials can contribute to the fixed output.
Jacobi proves $[f,g]\cdot F=f\cdot(g\cdot F)-g\cdot(f\cdot F)$.

For $p\geq0$, let $C^p$ be the space of jointly continuous alternating
$k$-multilinear maps $\mathfrak h^p\to A$. For $p=0$ this means $A$.
An alternating scalar map changes sign when two neighboring arguments
are exchanged and vanishes when they agree.

\begin{lemma}\label{lem:ham-core-finite-ghost-support}
Every jointly continuous scalar multilinear form on a finite power
of $\mathfrak h$ factors through a finite coefficient quotient in
each argument. Consequently
\begin{equation}\label{eq:ham-core-cochains}
 C^p=B^p:=\prod_\alpha u^\alpha\Lambda^pD.
\end{equation}
\end{lemma}

\begin{proof}
Continuity at zero gives open subspaces in the arguments on whose
product the form vanishes. Split each argument into its open subspace
and a finite-dimensional complement. The term using only open
subspaces vanishes. Every other term fixes at least one argument in
a finite-dimensional complement. Expand that argument in a finite
basis and apply induction to the resulting continuous forms with
fewer arguments. There are finitely many such terms, so their finite
quotients admit a common refinement. This proves the first assertion.

The exterior power $\Lambda^pD$ is the vector space of finite linear
combinations of alternating products of $p$ elements of $D$.
The first assertion identifies it with the continuous alternating
scalar forms. A map to the product $A$ is continuous exactly when
each coefficient map is continuous. Applying the scalar assertion
to each coefficient gives \eqref{eq:ham-core-cochains}.
\end{proof}

Write $c^I$ for the odd symbol corresponding to $\epsilon^I$.
The exterior product imposes $c^Ic^J=-c^Jc^I$ and $(c^I)^2=0$.
Thus an element of $B^p$ is
\[
 F=\sum_\alpha u^\alpha F_\alpha(c),
\]
where each $F_\alpha$ is a finite exterior polynomial of degree $p$.
Its set of ghost labels may depend on $\alpha$. This is the meaning
of a ghost coordinate here: an odd coordinate encoding an
antisymmetric symmetry parameter.

Give each $\Lambda^pD$ the discrete topology and $B^p$ the product
topology. Put $B=\bigoplus_{p\geq0}B^p$ as a graded algebra.
Completion takes place separately in each degree. Elements of the
underlying graded algebra have finitely many homogeneous components.
There is no imposed uniform ghost-support bound across the coefficients
of one component.

The product of two cochains is their alternating product. For degrees
$p,q$, it sums over permutations preserving the order of the first
$p$ and the last $q$ arguments, with the sign of that permutation.
In coordinates it is the product of the even $u_I$ and odd $c^I$.
Every output coefficient is a finite sum of products of finite exterior
polynomials. The product is therefore jointly continuous degreewise.

\section{The differential expressing the symmetry relations}

A coefficient $F\in A$ is invariant precisely when $f\cdot F=0$
for every $f\in\mathfrak h$. This is the first equation in a complex.
For $F\in C^p$, define
\begin{align*}
 (dF)(f_0,\ldots,f_p)
 &=\sum_{i=0}^p(-1)^i f_i\cdot
       F(f_0,\ldots,\widehat f_i,\ldots,f_p)\\
 &\quad+\sum_{0\leq i<j\leq p}(-1)^{i+j}
       F([f_i,f_j],f_0,\ldots,\widehat f_i,\ldots,
                                   \widehat f_j,\ldots,f_p).
\end{align*}
A hat means that the indicated argument is omitted; the remaining
arguments after the bracket retain their order. This differential
is called the Chevalley--Eilenberg differential. It combines the
action on coefficients with the bracket between symmetry parameters.

Write $[e_I,e_J]=\sum_K C^K_{IJ}e_K$. For
$I=(a,b)$ and $J=(c,d)$, its nonconstant part is
\begin{equation}\label{eq:ham-core-structure-constants}
 [e_I,e_J]=(ad-bc)x^{a+c-1}y^{b+d-1}\pmod{k}.
\end{equation}
A vanishing coefficient gives zero; a constant output is removed.
Evaluation of the differential on the coordinate generators gives
\begin{equation}\label{eq:ham-core-ce-differential}
 dc^K=-\frac12\sum_{I,J}C^K_{IJ}c^Ic^J,
 \qquad
 du_I=-\sum_{J,K}C^K_{IJ}u_Kc^J.
\end{equation}
The first sum is finite for each $K$. In the second, the coefficient
of a fixed $u_K$ has finite ghost support. Thus both expressions
belong to $B$.

The differential obeys $d(FG)=(dF)G+(-1)^pF(dG)$ for $F\in B^p$.
In the alternating product, an action term differentiates either
coefficient factor. A bracket term placing both bracketed arguments
in one factor gives its differential. Terms that split the bracketed
arguments between the two factors cancel in pairs, by exchanging
those arguments. Moving the differential past the first $p$ arguments
gives the displayed sign. This proves the product rule directly.

The coordinate formulas also prove continuity as a map $B^p\to B^{p+1}$.
For an output $u^\alpha$, a coefficient replacement uses a retained
label $K$ and an input label of degree at most $|K|+1$.
There are finitely many possible predecessor monomials. The ghost
differential acts on their discrete exterior-polynomial coefficients.
No bound uniform over all output monomials is needed.

\begin{proposition}\label{prop:ham-core-ce-square}
The operator $d$ has square zero. Thus $(B,d)$ is a complete
degreewise differential graded commutative algebra.
\end{proposition}

\begin{proof}
For a coefficient $u_I$, evaluation on $f,g$ gives
\[
 (d^2u_I)(f,g)=f\cdot(g\cdot u_I)-g\cdot(f\cdot u_I)
                                      -[f,g]\cdot u_I=0.
\]
For a scalar coordinate $c^K$, evaluation on $f,g,h$ gives
\[
 (d^2c^K)(f,g,h)
 =c^K\bigl([[f,g],h]+[[g,h],f]+[[h,f],g]\bigr)=0.
\]
The two equalities use the coefficient representation and Jacobi,
respectively. Applying the product rule twice cancels its two cross
terms, so $d^2(FG)=(d^2F)G+F(d^2G)$. It follows that $d^2$ vanishes
on finite coordinate polynomials. They are dense in each $B^p$.
Continuity of $d^2$ and separation of the coefficient topology extend
the equality to every cochain.
\end{proof}

\section{Cotangent coordinates and their degrees}

The coefficient action has a dual description. A continuous covector
$\lambda\in D$ transforms by the coadjoint action
\[
 (f\cdot\lambda)(g)=-\lambda([f,g]).
\]
If $\lambda$ uses coefficients through degree $N$, its transform uses
coefficients through degree $N+1$ and depends on $f$ through that degree.
The action into discrete $D$ is jointly continuous: near a fixed
$\lambda$, the second input can be fixed exactly. Jacobi proves the
representation identity for this action.

Use cochain shifts $(V[1])^n=V^{n+1}$. Place $\mathfrak h$ in degree
zero and $D[1]$ in degree minus one. Define the graded Lie algebra
\[
 G=\mathfrak h\ltimes D[1]
\]
by the bracket on $\mathfrak h$, the coadjoint action between the two
summands, and zero bracket between two elements of $D[1]$.
Graded antisymmetry means $[v,w]=-(-1)^{|v||w|}[w,v]$.
Its Jacobi identities are the Jacobi identity of $\mathfrak h$ and
the representation identity just proved; all cases with two covectors
vanish. The differential of $G$ is zero.

For completeness, the cochain sign convention can be described on
symmetric tensors of $G[1]$. Write $sv$ for the shifted element, of
degree $|v|-1$. The coaugmented symmetric coalgebra is
$\bigoplus_{n\geq0}S^n(G[1])$, including $S^0=k$.
Its coproduct is given by all ways of dividing
the inputs into two ordered groups. An exchange of homogeneous inputs
contributes its Koszul sign. The bracket extends to a coderivation
whose two-input part is
\[
 q_2(sv\,sw)=(-1)^{|v|}s[v,w].
\]
On more inputs, apply this formula to each pair and leave the other
inputs in their original order, with the sign from bringing the pair
together. Its square vanishes: disjoint pairs cancel, and the terms
using three inputs sum to graded Jacobi.

A homogeneous dual cochain $\phi$ has differential
$d\phi=-(-1)^{|\phi|}\phi q_2$. Products of dual cochains use the
coproduct and the Koszul sign for moving the second functional past
the first input group. For example, on two degree-zero Lie inputs,
\[
 (c^Ic^J)(se_A\,se_B)
 =-(\delta^I_A\delta^J_B-\delta^I_B\delta^J_A).
\]
The ordinary alternating form associated with an arity-$p$ functional
is $(-1)^{p(p-1)/2}\phi(sf_1\cdots sf_p)$.
These signs give exactly \eqref{eq:ham-core-ce-differential}.

The coordinate dual to $se_I$ is $c^I$ of standard cochain degree one.
The coordinate dual to $s^2\epsilon^I$ is $u_I$ of degree two.
Assign $u_I$ coefficient weight one and $c^I$ weight zero.
Subtract twice the weight from the standard cochain degree, then
complete over coefficient weights in each resulting degree.
The resulting degree-$p$ space is
\[
 \prod_{r\geq0}C^{p+2r}_{\mathrm{CE},\mathrm{cont}}(G)_{(r)}.
\]
The subscript $(r)$ specifies $r$ covector inputs. The remaining
$p$ inputs lie in $\mathfrak h$.

Because $D$ is discrete, continuity requires finite ghost support
for each fixed tuple of covectors. It imposes no common support bound
on all covector tuples. Symmetric multilinear forms in the covectors
are equivalent to their polynomial coefficients: evaluation on repeated
basis vectors contributes the corresponding factorials, which are
invertible in $k$. Lemma~\ref{lem:ham-core-finite-ghost-support}
therefore identifies the displayed product with $B^p$.
The negative coadjoint action dualizes to the adjoint coefficient
action, so this identification respects the differential as well as
the product. It proves an actual cotangent-cochain description of $B$.

\section{The contraction bracket}

A continuous $p$-fold multiderivation of $A$ is a jointly continuous
alternating map $A^p\to A$ satisfying the product rule in each argument.
For $p=0$ it is an element of $A$. Such maps are called continuous
polyvector fields. The coordinate vector field $\theta^I$ acts as
$\partial_{u_I}$; give it degree one.

Restriction of a multiderivation to the linear subspace
$\mathfrak h\subset A$ gives, at each output coefficient, a continuous
alternating scalar form. It therefore has finite ghost support by
Lemma~\ref{lem:ham-core-finite-ghost-support}. Conversely, any such
family of forms extends by the product rule. Each output coefficient
uses finitely many decompositions of its monomial and finitely many
derivative indices from those forms. This proves joint continuity in
the $p$ arguments. Density of finite coordinate polynomials proves
uniqueness. Consequently continuous polyvector fields are exactly $B^p$
under $c^I\mapsto\theta^I$. Give them the coefficient topology
\eqref{eq:ham-core-cochains}; this specifies the topology on the space
of operators, in addition to continuity of each operator.

There is a bracket of degree minus one on these polyvector fields,
the Schouten bracket. Its convention is fixed by
\[
 [\theta^I,u_J]_S=\delta^I_J,\qquad
 [\theta^I,\theta^J]_S=[u_I,u_J]_S=0.
\]
For $F$ of degree $p$, its full formula is
\begin{equation}\label{eq:ham-core-schouten}
 [F,G]_S=\sum_I\left(
 (-1)^{p+1}(\partial_{c^I}F)(\partial_{u_I}G)
          -(\partial_{u_I}F)(\partial_{c^I}G)\right).
\end{equation}
The exterior derivative is a left derivative:
$\partial_{c^I}(c^J F)=\delta_I^J F-c^J\partial_{c^I}F$.
On finite polynomials, the formula has graded antisymmetry
$[F,G]_S=-(-1)^{(p-1)(q-1)}[G,F]_S$ and the product rule of a
degree-minus-one bracket. Its Jacobiator is a derivation in each
argument: expanding a product cancels the cross terms in pairs.
It vanishes on triples of coordinate generators, because every inner
bracket there is constant. Induction on the three polynomial lengths
proves Jacobi for finite polynomials.

\begin{proposition}\label{prop:ham-core-schouten-continuity}
The bracket \eqref{eq:ham-core-schouten} is defined on all of $B$
and is jointly continuous degreewise. Its graded antisymmetry,
product rule, and Jacobi identity remain valid there.
\end{proposition}

\begin{proof}
Fix an output $u^\alpha$. A summand differentiates one factor in
$u_I$ and the other in $c^I$. After differentiation their coefficient
exponents satisfy $\beta+\gamma=\alpha$, with finitely many possible
pairs. The factor differentiated in $c^I$ is a finite exterior
polynomial. Its support permits only finitely many $I$. Thus the
coefficient is a finite sum and has finite ghost support.

For continuity at $(F,G)$, write the perturbations as $(U,V)$.
Require their coefficients at every divisor of $u^\alpha$ to vanish.
Then $[U,V]_S$ has zero coefficient at $u^\alpha$, because the
factor differentiated in a ghost direction still has one of those
coefficient monomials. The linear terms $[F,V]_S$ and $[U,G]_S$
use only finitely many additional monomials
$\beta+\mathbf e_I$, selected by the finite ghost supports of
the fixed coefficients of $F,G$. Fixing those perturbation
coefficients to zero kills the linear terms too. These conditions
define open neighborhoods. Repeating them for finitely many outputs
proves joint continuity. The polynomial identities extend by
continuity and density.
\end{proof}

Define the continuous bivector
\begin{equation}\label{eq:ham-core-pi}
 \Pi=\frac12\sum_{I,J,K}C^K_{IJ}u_K\theta^I\theta^J.
\end{equation}
For each nonconstant output $K$, the relation
$|I|+|J|=|K|+2$ makes its ghost sum finite. Hence $\Pi\in B^2$.
The elementary contraction is
\[
 [\theta^I\theta^J,u_K]_S
       =\delta^J_K\theta^I-\delta^I_K\theta^J.
\]
It gives $[\Pi,u_I]_S=du_I$ and $[\Pi,c^K]_S=dc^K$ with
\eqref{eq:ham-core-ce-differential}. Both sides are continuous
derivations, so they agree on all of $B$.
Every coefficient of $[\Pi,\Pi]_S$ is a finite Jacobi sum of the
structure constants \eqref{eq:ham-core-structure-constants}. Thus
\[
 [\Pi,\Pi]_S=0,\qquad
 d=[\Pi,-]_S,\qquad
 [\Pi,[\Pi,F]_S]_S=\tfrac12[[\Pi,\Pi]_S,F]_S=0.
\]
Together these formulas identify the cotangent cochains, the
coefficient-valued symmetry cochains, and the continuous polyvector
complex as differential graded algebras, with their specified
contraction bracket. In the standard cotangent cochain grading,
that contraction has degree minus three and weight minus one.
The regrading changes its degree to $-3-2(-1)=-1$.

For a formal Hamiltonian $f$, put $X_f=[\Pi,O_f]_S$.
The sign of its induced action is explicit:
\[
 [X_f,O_g]_S=-O_{[f,g]}.
\]
Thus the adjoint left action on coefficients is $-X_f$.
The entire translation action participates in this identity.

\section{An inverse limit that preserves the differential}

The preceding construction also produces an inverse limit of actual
cochain algebras. For $N\geq1$, let $J_N\subset B$ consist of
cochains whose coefficients vanish on all monomials of polynomial
degree at most $N$ with labels of Hamiltonian degree at most $N$.
Every ghost coordinate is retained. Each $J_N$ is a degreewise open
algebra ideal. Define
\[
 K_N=J_N\cap d^{-1}J_N,\qquad B_N=B/K_N.
\]
The extra condition says that both a cochain and its differential
vanish in the chosen coefficient window.

\begin{proposition}\label{prop:ham-core-cochain-tower}
Each $K_N$ is a differential graded ideal, and
\[
 J_{N+1}\subset K_N\subset J_N,
 \qquad
 B\xrightarrow{\ \sim\ }\varprojlim_N B_N
\]
is an isomorphism of complete degreewise differential graded
commutative algebras. The transition maps are surjective.
The quotients are infinite-dimensional, and their projections
do not preserve the Schouten bracket.
\end{proposition}

\begin{proof}
If $F\in K_N$, then $F,dF\in J_N$. The derivation rule implies
$FG,d(FG)\in J_N$ for every $G$, and $d^2F=0$ implies
$dF\in K_N$. This proves the ideal and differential assertions.
The differential preserves polynomial degree in the $u$ variables.
A retained output label $K$ can use only predecessor labels of
degree at most $|K|+1\leq N+1$. The ghost differential leaves
the coefficient monomial unchanged. Hence $dJ_{N+1}\subset J_N$,
which proves the two inclusions.

The $J_N$ form a cofinal system of neighborhoods for the coefficient
topology, and their quotient limit is the product defining $B$.
The inclusions show that the $K_N$ give the same topology and limit.
They are descending, so the quotient transition maps are surjective.
The maps respect the actual differential by construction.

All pure ghost polynomials survive in $B_N$; in particular its
degree-one part contains the infinite-dimensional space $D$.
If $|I|>N+1$, then $u_I\in J_{N+1}\subset K_N$, whereas
$[c^I,u_I]_S=1$. This excludes a bracket on $B_N$ compatible
with its projection.
\end{proof}

The bracket obstruction applies to every open quotient. A Schouten
ideal is an algebra ideal closed under brackets with every element
of $B$. The degree-zero part of an open ideal contains some
$J_N\cap B^0$, and therefore contains a sufficiently high $u_I$.
If the ideal is a Schouten ideal, contraction with $c^I$ places $1$
in it. Thus no proper open Schouten ideal exists.
The inverse limit preserves the differential and multiplication;
its individual quotients cannot retain all contractions.

\section{Changing the completion changes the operations}

There is a larger coefficient algebra
\[
 E^p=\prod_{\alpha,\,|S|=p}k u^\alpha c^S,
 \qquad E=\bigoplus_{p\geq0}E^p.
\]
Here $S$ is a finite set of ghost labels and $c^S$ is their ordered
exterior product. Choose once a total order on $\mathcal I$ to fix
that notation. In $E^p$, each individual ghost-monomial coefficient
has the discrete topology, and the whole space has the product topology.
There is a continuous injective algebra map $B\to E$.
Its image contains every finite polynomial and is therefore dense.
It is proper: $\sum_I c^I\in E^1$ has no representative in $B^1$.

The differential extends to $E$. A prescribed output word has only
finitely many differential predecessors. A ghost replacement uses
the bracket of two prescribed monomials. A coefficient replacement
uses the finite set of pairs producing a prescribed output label.
Thus every output coefficient of $dF$ remains a finite sum.
The inclusion is a map of differential graded algebras.

Its topology is not the topology induced from $E$. Choose labels
$I_n$ of strictly increasing degree. In $E$ both $u_{I_n}$ and
$c^{I_n}$ tend to zero, while
\[
 [c^{I_n},u_{I_n}]_S=1.
\]
The bracket therefore has no jointly continuous extension to $E$.
In $B$, the sequence $c^{I_n}$ does not tend to zero: its coefficient
at $u^0$ lies in the discrete space $D$. This is precisely the support
distinction used in the continuity proof.

The different symmetry sectors have correctly typed comparison
maps as well. Let $A_+$ be the coefficient algebra using only labels
of degree at least two. The adjoint action of $\mathfrak h_+$
preserves it. Restriction of Lie arguments and inclusion of coefficients
give continuous cochain algebra maps
\[
 C^\bullet_{\mathrm{CE},\mathrm{cont}}(\mathfrak h,A)
 \longrightarrow
 C^\bullet_{\mathrm{CE},\mathrm{cont}}(\mathfrak h_+,A)
 \longleftarrow
 C^\bullet_{\mathrm{CE},\mathrm{cont}}(\mathfrak h_+,A_+).
\]
Both maps have the middle complex as their target.

Even a Lie quotient does not automatically give a same-direction
cotangent map. Consider
\[
 q:\mathfrak m^2/\mathfrak m^4\longrightarrow
                    \mathfrak m^2/\mathfrak m^3.
\]
Let $r$ restrict covectors to the quadratic subspace. Take
$a=x^3$, $b=xy$, and $\lambda=(x^3)^\vee$. Then
\[
 r(a\cdot\lambda)(b)=-\lambda([a,b])=-3,
 \qquad ((qa)\cdot r\lambda)(b)=0.
\]
Thus $(q,r)$ does not respect the coadjoint action. The canonical
dual of $q$ is the opposite-direction injection $q^\vee$.
Additional cotangent transition data must satisfy their own
representation and contraction equations.

\section{Operators, tensor kernels, and matter representations}

An element of $B^1$ restricts to a continuous linear map
$\mathfrak h\to A$. The vector field
\[
 \mathcal E=\sum_Iu_Ic^I
\]
is permitted: the coefficient at each $u_I$ contains only $c^I$.
It acts on $A$ by polynomial degree and is the identity on the
linear subspace $\mathfrak h\subset A$.

The identity on $\mathfrak h$ is not an element of the algebraic
tensor product $\mathfrak h\otimes_kD$ under its map to operators,
\[
 \sum_{j=1}^r f_j\otimes\lambda_j
 \longmapsto\left(g\longmapsto
                  \sum_{j=1}^r\lambda_j(g)f_j\right).
\]
Every such operator has finite-dimensional image. The identity does
not. Allowing coefficientwise finite ghost support produces a larger
operator space. A tensor completion intended to represent these
operators must be specified and compared with that space; the
existence of $\mathcal E$ alone does not make the identity a kernel
in every chosen tensor product.

There is also a direct restriction on finite-dimensional matter.
For a finite-dimensional vector space $V$ with discrete topology,
a continuous representation is a continuous Lie map
$\rho:\mathfrak h\to\operatorname{End}_k(V)$, where the target
has the discrete topology and commutator bracket.
Proposition~\ref{prop:ham-core-no-finite-quotient}
forces $\rho=0$. Consequently the finite-dimensional current coupling
of Chapter~\ref{chap:quantum-current-anomaly} has no nontrivial
instance for the full formal Hamiltonian algebra with this topology.
An origin-preserving symmetry algebra, an infinite-dimensional matter
space, or a changed coefficient topology gives a different problem.
The latter choices require their own continuous action, tensor
contractions, and trace or regularization. These constructions are
the next step toward a quantum realization of the full Hamiltonian
symmetry.
